\documentclass[11pt]{article}
\usepackage[margin=25mm]{geometry}
\usepackage{fontspec}
\setmainfont{Latin Modern Roman}
\usepackage{amsmath,amssymb,mathtools}
\usepackage{unicode-math}
\setmathfont{Latin Modern Math}
\usepackage{xurl,hyperref}
\hypersetup{colorlinks=true,urlcolor=blue}
\setcounter{secnumdepth}{0}
\setlength{\parindent}{0pt}
\setlength{\parskip}{5pt}
\setlength{\emergencystretch}{3em}
\newcommand{\sref}[2]{\hyperlink{source:#1:#2}{[S#2]}}
\newcommand{\cataloguescope}[1]{\paragraph{Recorded target.}#1}
\begin{document}
\section{A d-dimensional permutation of order n is an n-by-n-by}\label{statement}

\subsection{Definitions and mathematical statement}

\paragraph{Shared notation.}$\mathbb N=\{1,2,\ldots\}$, and $\mathbb Z,\mathbb Q,\mathbb R,\mathbb C$ denote the integers, rationals, reals and complex numbers. Any other conventions are specified locally.


Fix an integer $d\geq1$ and write $[n]=\{1,\ldots,n\}$. A $d$-dimensional permutation of order $n$ is a function $A:[n]^{d+1}\to\{0,1\}$ such that for each $k\in\{1,\ldots,d+1\}$ and each fixed choice of the other $d$ coordinates,
\[\sum_{t=1}^n A(j_1,\ldots,j_{k-1},t,j_{k+1},\ldots,j_{d+1})=1.\]
Let $P(n,d)$ count these arrays, so $P(n,1)=n!$ and $P(n,2)$ counts Latin squares of order $n$. The original source establishes an upper bound with the asymptotic scale below.
\paragraph{Statement recorded in the supplied source.}
For every fixed $d\geq1$, does the matching lower bound hold as $n\to\infty$, and hence
\[P(n,d)=\left((1+o(1))\frac{n}{e^d}\right)^{n^d}?\]
Equivalently, is $\log P(n,d)=n^d(\log n-d+o(1))$ for each fixed $d$? \sref{statement}{2}
\subsection{Short English statement}
Does the number of fixed-dimensional permutation arrays have the asymptotic size predicted by the Linial--Luria upper bound?
\subsection{Sources}
\begin{description}
\item[\hypertarget{source:statement:1}{S1}] ULAM.AI and the UnsolvedMath Contributors, UnsolvedMath: A Curated Collection of Open Mathematics Problems, catalogue entry. CC BY 4.0. \url{https://huggingface.co/datasets/ulamai/UnsolvedMath}.
\item[\hypertarget{source:statement:2}{S2}] Nathan Linial and Zur Luria, An Upper Bound on the Number of High-dimensional Permutations (2012), Definition 1.2, Conjecture 1.3 and Theorem 1.4, pp. 2--4. \url{https://arxiv.org/abs/1106.0649}.
\end{description}
\label{end:statement}
\end{document}
